% Part: incompleteness
% Chapter: representability-in-q
% Section: representing-relations

\documentclass[../../../include/open-logic-section]{subfiles}

\begin{document}

\olfileid[hi]{inc}{req}{rel}

\olsection{संबंधों का निरूपण}

आइए बताएँ कि किसी \emph{संबंध} के निरूपणीय होने का क्या अर्थ है।

\begin{defn}
\ollabel{defn:representing-relations} प्राकृतिक संख्याओं पर संबंध
$R(x_0,\dots,x_k)$ {\em $\Th{Q}$ में निरूपणीय} है यदि कोई सूत्र
$!A_R(x_0,\dots,x_k)$ ऐसा हो कि जब भी $R(n_0,\dots,n_k)$ सत्य हो,
$\Th{Q}$,~$!A_R(\num{n_0},\dots,\num{n_k})$ को सिद्ध करे, और जब भी
$R(n_0,\dots,n_k)$ असत्य हो, $\Th{Q}$,~$\lnot !A_R(\num{n_0},
\dots, \num{n_k})$ को सिद्ध करे।
\end{defn}

\begin{thm}
\ollabel{thm:representing-rels} कोई संबंध $\Th{Q}$ में तभी और केवल तभी
निरूपणीय है जब वह संगणनीय है।
\end{thm}

\begin{proof}
आगे की दिशा के लिए मान लें कि सूत्र $!A_R(x_0,\dots,x_k)$ संबंध
$R(x_0,\dots,x_k)$ का निरूपण करता है। $R$ को संगणित करने की एक
कलन-विधि यह है: निवेश $n_0$, \dots,~$n_k$ पर एक साथ
$!A_R(\num{n_0}, \dots, \num{n_k})$ के प्रमाण और
$\lnot !A_R(\num{n_0}, \dots, \num{n_k})$ के प्रमाण की खोज करें। हमारी
मान्यता के अनुसार खोज को दोनों में से कोई एक अवश्य मिलेगा; यदि पहला मिले,
तो ``हाँ'' बताएँ, अन्यथा ``नहीं''।

दूसरी दिशा में मान लें कि $R(x_0, \dots, x_k)$ संगणनीय है। परिभाषा के
अनुसार इसका अर्थ है कि फलन $\Char{R}(x_0, \dots, x_k)$ संगणनीय है।
\olref[int]{thm:representable-iff-comp} के अनुसार $\Char{R}$ का निरूपण
किसी सूत्र, मान लें $!A_{\Char{R}}(x_0, \dots, x_k, y)$, द्वारा होता है।
$!A_R(x_0, \dots, x_k)$ को सूत्र
$!A_{\Char{R}}(x_0, \dots, x_k, \num{1})$ मानें। तब किन्हीं भी
$n_0$, \dots,~$n_k$ के लिए, यदि $R(n_0, \dots, n_k)$ सत्य है, तो
$\Char{R}(n_0, \dots, n_k) = 1$; उस स्थिति में $\Th{Q}$,
$!A_{\Char{R}}(\num{n_0}, \dots, \num{n_k}, \num{1})$ को सिद्ध करता है,
और इसलिए $\Th{Q}$,~$!A_R(\num{n_0}, \dots, \num{n_k})$ को सिद्ध करता है।
दूसरी ओर यदि $R(n_0, \dots, n_k)$ असत्य है, तो
$\Char{R}(n_0, \dots, n_k) = 0$। इसका अर्थ है कि $\Th{Q}$ सिद्ध करता है
\[
\lforall[y][(!A_{\Char{R}}(\num{n_0}, \dots, \num{n_k}, y) \lif y =
  \num{0})].
\]
चूँकि $\Th{Q}$,~$\eq/[\num{0}][\num{1}]$ को सिद्ध करता है, इसलिए
$\Th{Q}$,~$\lnot !A_{\Char{R}}(\num{n_0}, \dots, \num{n_k}, \num{1})$
को सिद्ध करता है और इस प्रकार
$\lnot !A_R(\num{n_0}, \dots, \num{n_k})$ को भी सिद्ध करता है।
\end{proof}

\begin{prob}
दिखाएँ कि यदि $R$,~$\Th{Q}$ में निरूपणीय है, तो $\Char{R}$ भी निरूपणीय है।
\end{prob}

\end{document}
